% ============================================================
% SPUN-OUT MATERIAL (removed from optimal_long_short.tex per revision task T2).
% This "barbell" section was cut from the main manuscript because the keeper leg
% is not an investable position (liquidation is a latency/MEV/gas race), the
% weights were normalised per unit of another agent's equity, and the backtest
% was under-specified and underperformed plain USDC supply. Preserved here
% verbatim so nothing is lost; the Jensen-gap observation (eq:jensen) was kept
% in the main text as a remark in the regime-switching subsection.
% ============================================================

% ============================================================
\section{A Barbell Long the Two Convexities}\label{sec:barbell}
% ============================================================

The preceding sections assemble, without setting out to, a complete indictment of the leveraged
long--short book. Section~\ref{sec:carry-crosssection} shows it collects no premium: every
volatile collateral pair on Aave carries \emph{negative} net carry, from $-1.8\%$ to $-4.5\%$ per
year, because the wrapped-BTC and staked-ETH reserves pay almost nothing as collateral while their
debt legs cost several percent. Yet the same book is \emph{short two convex payoffs}. It is short
the barrier: the killed payoff $\Pi_T$ forfeits the recovery $R_\tau$ of Section~\ref{sec:recovery}
at liquidation, a bounded upside sold against a rare near-total loss. And it is short the
utilization kink: the borrower pays the piecewise-linear rate $r_b(U)$ of
Section~\ref{sec:carry-regime}, whose slope jumps by more than a factor of $25$ at $U_{\mathrm{opt}}$
(published $\mathrm{slope}_2$ is $0.60$--$0.80$ for WETH and the stablecoins against a
first-segment slope an order of magnitude smaller, and $3.0$ for the BTC wrappers), so the borrower
is short a call on utilization. A position that pays a negative premium to stay short two
convexities is fragile in exactly Taleb's sense.

The constructive response is to take the other side of both, and both sides are available
\emph{on-protocol}, with no options market. We place a large allocation $w_{\mathrm{safe}}$ in a
pure supplier so deeply buffered that the barrier is unreachable --- the $h_0\to\infty$ limit of
Remark~\ref{rem:no-shorting}, in which the payoff converges to a unit long and $p_{\mathrm{surv}}\to1$
--- and a small allocation $w_{\mathrm{keep}}=1-w_{\mathrm{safe}}$ that buys the two convexities the
levered book was short. The construction reuses the machinery already built; it introduces no new
transforms.

\subsection{The keeper leg: long the barrier}\label{sec:barbell-keeper}
A liquidator repays a fraction $\phi_c$ of the outstanding debt (Aave's close factor, $0.5$ at the
boundary) and seizes collateral worth $(1+\lambda)$ times the repayment, so on $\{\tau\le T\}$ the
keeper collects $\lambda\phi_c V_{2,\tau}$. Normalised per unit of the borrower's initial equity,
$n_2 S_{2,0}=b/(e^{h_0}-b)$, its expected value is
\begin{equation}\label{eq:keeper}
\mathcal K(h_0)\;=\;\lambda\,\phi_c\,\frac{b}{e^{h_0}-b}\;\E\!\left[e^{X_{2,\tau}}\one_{\{\tau\le T\}}\right],
\end{equation}
which is \eqref{eq:fp-transform} at $k=1$ --- the \emph{same} first-passage transform that priced the
recovery $R_\tau$ in Section~\ref{sec:recovery}, evaluated on the same $3\times3$ barrier system.
What the borrower forfeits at liquidation and what the keeper collects are one object read with
opposite signs: $R_\tau=V_{2,\tau}(1-b(1+\lambda))/b$ against $\lambda\phi_c V_{2,\tau}$, both
proportional to $\E[e^{X_{2,\tau}}\one_{\{\tau\le T\}}]$. The keeper leg is therefore priced exactly
and for free by the transform already validated to $1.7\times10^{-11}$ in
Section~\ref{sec:recovery}; being long it is being long precisely the liquidation event that kills
the borrower.

\subsection{The supplier leg: long the kink}\label{sec:barbell-supplier}
A pure supplier bears no barrier and earns $(1-\phi)\,U\,r_b(U)$, with $\phi$ the reserve factor.
Because $r_b$ is convex in $U$ --- linear with an upward slope break at $U_{\mathrm{opt}}$ --- so is
the supplier's rate, and Jensen's inequality makes the value of the kink explicit. Modelling
utilization as a Jacobi (Wright--Fisher) diffusion
$dU_t=\kappa(\theta-U_t)\,dt+\sigma\sqrt{U_t(1-U_t)}\,dW_t$, which unlike a CIR fit respects
$U\le1$, its stationary law is $\mathrm{Beta}(a,b)$ with $a=2\kappa\theta/\sigma^2$,
$b=2\kappa(1-\theta)/\sigma^2$, and the supplier's expected yield exceeds the yield evaluated at
mean utilization by the strictly positive gap
\begin{equation}\label{eq:jensen}
\mathcal J\;=\;\underbrace{\E\!\big[(1-\phi)\,U\,r_b(U)\big]}_{\text{expected yield}}
\;-\;\underbrace{(1-\phi)\,\E[U]\,r_b(\E[U])}_{\text{yield at mean }U}\;>\;0 .
\end{equation}
The gap $\mathcal J$ is the premium for being long the kink, and it is invisible to any
deterministic-rate model --- including the carry substitution of Proposition~\ref{prop:carry},
which by construction evaluates the curve at a single $U$ and so prices $\mathcal J$ at zero. This is
the precise sense in which Section~\ref{sec:carry-crosssection}'s ``deterministic carry is inert''
is a statement about the \emph{first} moment only: the second-order curvature the borrower pays for,
the supplier is paid for, and the barbell moves it from the liability to the asset side. The two legs
are implemented in \texttt{optimal\_long\_short.barbell} (\texttt{keeper\_leg}, \texttt{supplier\_leg}),
the keeper leg calling the identical inverter as \texttt{recovery}.

\subsection{Empirical behaviour}\label{sec:barbell-empirical}
A historical backtest on Aave v3 Ethereum (Jan 2023--Jun 2026, $1{,}222$ daily observations, USDC
supply against WETH borrow) compares a $90/10$ barbell against the naive carry book and three
conventional alternatives. The naive carry trade is liquidated to a total loss over the window
(cumulative return $-100\%$, maximum drawdown $-100\%$); the barbell instead delivers a positive
cumulative $+15.1\%$ with its drawdown structurally bounded at $-2.6\%$, the best Calmar ratio in the
comparison ($1.66$). This is the qualitative signature the construction predicts: the safe leg floors
the loss while the two convex legs supply small, steady, positively skewed gains.

Two honest caveats attach, both consistent with the paper's thesis rather than against it. First,
the keeper leg's payoff inherits the continuous-crossing optimism of Remark~\ref{rem:overshoot}: a
jump that overshoots the barrier leaves the liquidated borrower with bad debt, and a keeper who is
not first in the queue collects less than \eqref{eq:keeper} returns, so $\mathcal K$ is an upper
bound. Second, on a \emph{risk-adjusted} basis the barbell does not clear the on-chain hurdle: over
this window the mean USDC supply rate was $4.80\%$, and no strategy in the comparison --- the barbell
included, at a $4.3\%$ annualised return --- earned a positive Sharpe against simply supplying USDC.
The barbell wins on tail geometry (bounded drawdown, positive Calmar, positive skew), not on excess
mean return. That is the correct reading of a convexity-harvesting trade in a regime where the
risk-free itself is a healthy on-chain yield: it is bought for what it does in the left tail, and it
is the disciplined answer to the fragility the rest of the paper diagnoses.
